\documentclass{article}
\usepackage[UTF8]{ctex}
\usepackage{amsmath,amssymb,amsfonts}
\usepackage{enumitem}
\usepackage{geometry}
\geometry{a4paper, margin=1in}
% === 错题本排版优化（不参与题目标识） ===
\setlength{\parindent}{0pt}
\setlength{\parskip}{0.35em plus 0.10em minus 0.08em}
\setlength{\emergencystretch}{3em}
\setlist{itemsep=0.25em, topsep=0.40em, parsep=0pt}
\allowdisplaybreaks[2]
% === 错题本排版优化结束 ===
\begin{document}
\section*{广西大学2024年数学分析真题}
\section*{162．广西大学2024年数学分析试卷（回忆版）}

\subsection{题目1}
1．求极限 \(\lim _{x \rightarrow 0} \frac{\sqrt{1+\tan x}-\sqrt{1-\sin x}}{x}\) ．

\subsection{题目2}
2．设 \(z=x^{2}+y^{2}\) ，其中 \(y=f(x)\) 是由方程 \(x^{2}-x y+y^{2}=1\) 确定的隐函数，求 \(\frac{\mathrm{d}^{2} z}{\mathrm{~d} x^{2}}\) ．

\subsection{题目3}
3．未知。

\subsection{题目4}
4．抛物线 \(y^{2}=2 x\) 将 \(x^{2}+y^{2} \leq 8\) 分成左右两部分，求左右两部分面积之比．

\subsection{题目5}
5．求 \(y^{2}=2 p x\) 上一点，使得此点处的法线被抛物线截取的线段最短．

\subsection{题目6}
6．求曲线积分 \(\int_{L} x y \mathrm{~d} s\) ，其中 \(L\) 为椭圆 \(\frac{x^{2}}{a^{2}}+\frac{y^{2}}{b^{2}}=1\) 在第一象限的部分。

\subsection{题目7}
7．求重积分 \(\iiint_{V}\left(x^{2}+y^{2}+z\right) \mathrm{d} x \mathrm{~d} y \mathrm{~d} z\) ，其中 \(V\) 是直线 \(z=y, x=0\) 绕 \(z\) 轴旋转所得曲面与 \(z=1\)围成的区域。

\subsection{题目8}
8．求 \(y=\arcsin (\cos x)\) 的傅里叶级数展开式。

\subsection{题目9}
9 ．求曲面积分 \(\iint_{S}(x+y+z) \mathrm{d} S\) ，其中 \(S: x^{2}+y^{2}+z^{2}=a^{2}, z \geq 0\) ．

\subsection{题目10}
10．已知 \(\lim _{n \rightarrow \infty} a_{n}=a\) ，若 \(a_{n}, a>0\) ，证明： \(\lim _{n \rightarrow \infty} \sqrt[n]{a_{n}}=1\) ．

\subsection{题目11}
11．设 \(f(x)\) 在 \([a,+\infty)\) 上连续，且 \(\lim _{x \rightarrow+\infty} f(x)\) 存在，证明：\(f(x)\) 在 \([a,+\infty)\) 上一致连续．

\subsection{题目12}
12．证明 \(f(x)=x^{3} e^{-x^{2}}\) 是有界函数．

\subsection{题目13}
13．证明 \(\int_{0}^{+\infty} \frac{\sin x y}{y} \mathrm{~d} y\) 在 \([\delta,+\infty)(\delta>0)\) 上一致收敛，在 \((0,+\infty)\) 上不一致收敛。

\subsection{题目14}
14．由闭域套定理证明 \(\mathbb{R}^{2}\) 中的有界无限点集 \(E\) 在 \(\mathbb{R}^{2}\) 中至少有一个聚点．

\subsection{题目15}
15．设 \(S\) 是包围 \(V\) 的光滑曲面，证明 \(V\) 的体积 \(\Delta V\) 为

\[
\Delta V=\frac{1}{3} \iint_{S}(x \cos \alpha+y \cos \beta+z \cos \gamma) \mathrm{d} S
\]

其中 \(\cos \alpha, \cos \beta, \cos \gamma\) 是曲面 \(S\) 外法线的方向余弦。
\end{document}
